%  MDPI Information, article. Built from template.tex of MDPI_template_ACS.zip (class files in Definitions/, unchanged).
%  Every number of the text is a \val{key} from numbers.tex (generated by tools/paper_numbers.py); see README.md.

%=================================================================
\documentclass[information,article,submit,moreauthors]{Definitions/mdpi}

%=================================================================
% MDPI internal commands - do not modify
\firstpage{1}
\makeatletter
\setcounter{page}{\@firstpage}
\makeatother
\pubvolume{1}
\issuenum{1}
\articlenumber{0}
\pubyear{2026}
\copyrightyear{2026}
\datereceived{ }
\daterevised{ }
\dateaccepted{ }
\datepublished{ }

%=================================================================
% No packages beyond those of the class (it loads graphicx, booktabs, tabularx, xcolor, amsmath, etoolbox, geometry,
% tikz for the ORCID icon, and others). Figure 1 is a precompiled PDF (assets/cascade_diagram.tex), so no TikZ
% libraries are needed here.
% Faster compilation: the class inserts the MDPI logo as Definitions/logo-mdpi.eps, which pdfLaTeX would convert
% with epstopdf on every run; the same logo, converted once with Ghostscript, is used as a PDF instead.
\makeatletter
\patchcmd{\@maketitle}{Definitions/logo-mdpi.eps}{Definitions/logo-mdpi.pdf}{}{\typeout{MDPI logo: patch not applied}}
\patchcmd{\@maketitle}{Definitions/logo-mdpi.eps}{Definitions/logo-mdpi.pdf}{}{}
\makeatother
% fancyhdr asks for a taller header box; with includehead, geometry keeps the page margins and shortens the body
\geometry{headheight=21pt}
% Generated by paper_mdpi/tools/paper_numbers.py from results/, reports/ and configs/. Do not edit by hand.
\makeatletter
\@namedef{val@men-raw}{5971}% reports/stats/dedup_mendeley.json rows_before
\@namedef{val@men-clean}{5739}% reports/stats/dedup_mendeley.json rows_after
\@namedef{val@men-in-uci}{4914}% reports/stats/overlap.json mendeley_uci[clean] exact_or_near
\@namedef{val@men-in-uci-pct}{85.6}% reports/stats/overlap.json / dedup_mendeley.json
\@namedef{val@men-removed}{232}% reports/stats/dedup_mendeley.json rows_before - rows_after
\@namedef{val@men-test-n}{1148}% reports/stats/split_mendeley.json
\@namedef{val@men-test-smishing}{99}% reports/stats/split_mendeley.json
\@namedef{val@men-test-spam}{84}% reports/stats/split_mendeley.json
\@namedef{val@shim-test-n}{183}% reports/stats/split_mendeley.json (smishing + spam)
\@namedef{val@ssc-removed}{4431}% reports/stats/super_sms_clean.json removed
\@namedef{val@men-conflict-groups}{41}% reports/stats/dedup_mendeley.json conflict_groups (exact + near)
\@namedef{val@cascade-clean-f1}{0.9316}% results/final/metrics.csv (mean over 5 seeds)
\@namedef{val@cascade-clean-f1-sd}{0.0074}% results/final/metrics.csv (mean over 5 seeds)
\@namedef{val@svmchar-clean-f1}{0.9374}% results/final/metrics.csv (mean over 5 seeds)
\@namedef{val@svmchar-clean-f1-sd}{0.0000}% results/final/metrics.csv (mean over 5 seeds)
\@namedef{val@svmcharaug-clean-f1}{0.9310}% results/final/metrics.csv (mean over 5 seeds)
\@namedef{val@svmcharaug-clean-f1-sd}{0.0036}% results/final/metrics.csv (mean over 5 seeds)
\@namedef{val@distilaug-clean-f1}{0.9149}% results/final/metrics.csv (mean over 5 seeds)
\@namedef{val@distilaug-clean-f1-sd}{0.0115}% results/final/metrics.csv (mean over 5 seeds)
\@namedef{val@bert-clean-f1}{0.9195}% results/final/metrics.csv (mean over 5 seeds)
\@namedef{val@bert-clean-f1-sd}{0.0039}% results/final/metrics.csv (mean over 5 seeds)
\@namedef{val@clean-best-p-boot-holm}{0.088}% results/final/significance.csv
\@namedef{val@clean-best-p-mcn-holm}{0.180}% results/final/significance.csv
\@namedef{val@routed-min-pct}{2.0}% results/final/metrics.csv routed_share (mean over seeds)
\@namedef{val@routed-max-pct}{2.8}% results/final/metrics.csv routed_share (mean over seeds)
\@namedef{val@routed-binary-min-pct}{45}% results/final/metrics.csv routed_share, Mahmud and transfer
\@namedef{val@routed-binary-max-pct}{77}% results/final/metrics.csv routed_share, Mahmud and transfer
\@namedef{val@pub-ho20-min}{0.4416}% results/final/metrics.csv (mean over 5 seeds)
\@namedef{val@pub-ho20-max}{0.8555}% results/final/metrics.csv (mean over 5 seeds)
\@namedef{val@cascade-mahmud-ho20}{0.9840}% results/final/metrics.csv (mean over 5 seeds)
\@namedef{val@svmchar-mahmud-ho20}{0.9870}% results/final/metrics.csv (mean over 5 seeds)
\@namedef{val@cascade-shim-ho20}{0.9080}% results/final/metrics.csv (mean over 5 seeds)
\@namedef{val@svmchar-shim-ho20}{0.9233}% results/final/metrics.csv (mean over 5 seeds)
\@namedef{val@lat-cascade-p50}{1.60}% results/final/cascade_end_to_end.json
\@namedef{val@lat-cascade-p95}{2.47}% results/final/cascade_end_to_end.json
\@namedef{val@lat-cascade-thr}{297}% results/final/cascade_end_to_end.json
\@namedef{val@lat-cascade-routed-pct}{2.4}% results/final/cascade_end_to_end.json
\@namedef{val@lat-distilaug-p50}{41.8}% results/final/latency.csv
\@namedef{val@lat-svmcharaug-p50}{1.46}% results/final/latency.csv
\@namedef{val@n-seeds}{5}% configs/final.yaml seeds
\@namedef{val@seed-last}{5}% configs/final.yaml seeds
\@namedef{val@shingle-k}{5}% src/data/common.py SHINGLE_K
\@namedef{val@jaccard-thr}{0.9}% src/data/common.py JACCARD_THRESHOLD
\@namedef{val@num-perm}{128}% src/data/common.py NUM_PERM
\@namedef{val@lsh-thr}{0.7}% src/data/common.py
\@namedef{val@check-lsh-thr}{0.5}% src/data/split.py
\@namedef{val@split-train}{70}% src/data/split.py FRACTIONS
\@namedef{val@split-val}{10}% src/data/split.py FRACTIONS
\@namedef{val@split-test}{20}% src/data/split.py FRACTIONS
\@namedef{val@split-seed}{42}% src/data/common.py SEED
\@namedef{val@ece-bins}{15}% configs/step2.yaml ece_bins
\@namedef{val@lvl-a}{5}% configs/step4c.yaml eval_levels
\@namedef{val@lvl-b}{10}% configs/step4c.yaml eval_levels
\@namedef{val@lvl-c}{20}% configs/step4c.yaml eval_levels
\@namedef{val@obf-seed}{2026}% src/eval/obfuscation.py SEED
\@namedef{val@aug-p}{0.5}% configs/step4c.yaml augmentation
\@namedef{val@aug-max}{0.2}% configs/step4c.yaml augmentation
\@namedef{val@char-ng-lo}{2}% configs/step4c.yaml char_branch
\@namedef{val@char-ng-hi}{5}% configs/step4c.yaml char_branch
\@namedef{val@svm-c}{0.583}% results/best_params/mendeley_3class/linear_svm.json
\@namedef{val@calib-folds}{5}% configs/step2.yaml linear_svm
\@namedef{val@distil-lr}{6.19\times 10^{-5}}% configs/step4a.yaml text.lr
\@namedef{val@distil-batch}{32}% configs/step4a.yaml text
\@namedef{val@distil-epochs}{5}% configs/step4a.yaml text
\@namedef{val@distil-patience}{2}% configs/step4a.yaml text
\@namedef{val@max-len}{128}% configs/step4a.yaml text
\@namedef{val@fusion-d}{256}% configs/step4a.yaml fusion.d
\@namedef{val@aux-w}{0.5}% configs/step4a.yaml fusion
\@namedef{val@hand-n}{14}% src/fusion/features.py FEATURES
\@namedef{val@svd-k}{256}% configs/step4c.yaml char_branch
\@namedef{val@tau}{0.6651}% configs/final.yaml method.tau
\@namedef{val@tau-noaug}{0.5772}% configs/final.yaml ablations
\@namedef{val@tau-seeds}{1--3}% scripts/freeze_tau.py: seeds with step-4c val probabilities
\@namedef{val@cost-r1}{1}% configs/final.yaml cost_ratios
\@namedef{val@cost-r2}{5}% configs/final.yaml cost_ratios
\@namedef{val@cost-r3}{10}% configs/final.yaml cost_ratios
\@namedef{val@boot-n}{10\,000}% configs/final.yaml significance
\@namedef{val@boot-seed}{42}% configs/final.yaml significance
\@namedef{val@alpha}{0.05}% configs/final.yaml significance
\@namedef{val@mcn-exact}{25}% src/eval/stats.py mcnemar
\@namedef{val@optuna-trials}{20}% configs/step2.yaml n_trials
\@namedef{val@qwen-shots}{8}% configs/step3.yaml qwen_fewshot
\@namedef{val@bert-epochs}{5}% configs/step3.yaml bert
\@namedef{val@ld-a}{1}% configs/final.yaml low_data fractions
\@namedef{val@ld-b}{5}% configs/final.yaml low_data fractions
\@namedef{val@ld-c}{10}% configs/final.yaml low_data fractions
\@namedef{val@ld-d}{25}% configs/final.yaml low_data fractions
\@namedef{val@ld-min}{5}% configs/step4b.yaml low_data
\@namedef{val@tr-frac}{10}% configs/step4b.yaml transfer
\@namedef{val@mh-w2v}{100}% configs/step5a.yaml mahmud
\@namedef{val@mh-filters}{32}% configs/step5a.yaml mahmud
\@namedef{val@mh-kernel}{3}% configs/step5a.yaml mahmud
\@namedef{val@mh-pool}{2}% configs/step5a.yaml mahmud
\@namedef{val@mh-gru}{64}% configs/step5a.yaml mahmud
\@namedef{val@mh-drop}{0.2}% configs/step5a.yaml mahmud
\@namedef{val@mh-batch}{128}% configs/step5a.yaml mahmud
\@namedef{val@mh-epochs}{50}% configs/step5a.yaml mahmud
\@namedef{val@mh-maxlen}{100}% configs/step5a.yaml mahmud
\@namedef{val@sh-batches}{8, 16, 32}% configs/step5a.yaml shim
\@namedef{val@sh-lrs}{1\times 10^{-6}, 5\times 10^{-6}, 1\times 10^{-5}, 5\times 10^{-5}}% configs/step5a.yaml shim
\@namedef{val@sh-epochs}{10}% configs/step5a.yaml shim
\@namedef{val@sh-inner}{20}% configs/step5a.yaml shim
\@namedef{val@lat-msgs}{1\,000}% configs/final.yaml latency
\@namedef{val@lat-warmup}{20}% configs/final.yaml latency
\@namedef{val@clean-diff-abs}{1.85}% results/final/significance.csv (seed 1, cascade below)
\@namedef{val@clean-ci-lo}{$-$4.12}% results/final/significance.csv
\@namedef{val@clean-ci-hi}{$+$0.32}% results/final/significance.csv
\@namedef{val@qwen-clean-f1}{0.8202}% results/final/metrics.csv (mean over 5 seeds)
\@namedef{val@cm-spam-to-smish}{8}% results/final/confusion_cascade_seed1.csv
\@namedef{val@cm-smish-to-spam}{9}% results/final/confusion_cascade_seed1.csv
\@namedef{val@cm-ham-wrong}{4}% results/final/confusion_cascade_seed1.csv
\@namedef{val@cm-ham-n}{965}% results/final/confusion_cascade_seed1.csv
\@namedef{val@ece-cascade}{0.0131}% results/final/metrics.csv (mean over 5 seeds)
\@namedef{val@ece-svmcharaug}{0.0184}% results/final/metrics.csv (mean over 5 seeds)
\@namedef{val@ece-distilaug}{0.0170}% results/final/metrics.csv (mean over 5 seeds)
\@namedef{val@ss-min}{0.9872}% results/final/metrics.csv (mean over 5 seeds), super_sms without the few-shot model
\@namedef{val@ss-max}{0.9936}% results/final/metrics.csv (mean over 5 seeds), super_sms without the few-shot model
\@namedef{val@ss-cascade}{0.9934}% results/final/metrics.csv (mean over 5 seeds)
\@namedef{val@ss-bert}{0.9936}% results/final/metrics.csv (mean over 5 seeds)
\@namedef{val@ss-p-boot}{0.214}% results/final/significance.csv
\@namedef{val@ss-p-mcn}{0.243}% results/final/significance.csv
\@namedef{val@claim-mh-acc}{0.9817}% results/step5a/published_claims.yaml (from the paper)
\@namedef{val@claim-sh-min}{91.1}% results/step5a/published_claims.yaml (from the paper, %)
\@namedef{val@claim-sh-max}{92.9}% results/step5a/published_claims.yaml (from the paper, %)
\@namedef{val@mh-ho20}{0.8555}% results/final/metrics.csv (mean over 5 seeds)
\@namedef{val@sh-ho20-min}{0.4416}% results/final/metrics.csv (mean over 5 seeds)
\@namedef{val@sh-ho20-max}{0.6847}% results/final/metrics.csv (mean over 5 seeds)
\@namedef{val@realised-ho20}{20.0}% results/final/levels/mendeley_3class_heldout20.json
\@namedef{val@rob-neural-ho20-min}{0.6073}% results/final/metrics.csv (mean over 5 seeds)
\@namedef{val@rob-neural-ho20-max}{0.6998}% results/final/metrics.csv (mean over 5 seeds)
\@namedef{val@rob-distilaug-ho20}{0.8872}% results/final/metrics.csv (mean over 5 seeds)
\@namedef{val@rob-distil-ho20}{0.6462}% results/final/metrics.csv (mean over 5 seeds)
\@namedef{val@rob-svmchar-ho20}{0.9322}% results/final/metrics.csv (mean over 5 seeds)
\@namedef{val@rob-cascade-ho20}{0.9326}% results/final/metrics.csv (mean over 5 seeds)
\@namedef{val@rob-v1-ho20}{0.6197}% results/final/metrics.csv (mean over 5 seeds)
\@namedef{val@rob-v2-ho20}{0.9232}% results/final/metrics.csv (mean over 5 seeds)
\@namedef{val@routed-clean-pct}{2.1}% results/final/metrics.csv (mean over 5 seeds) routed_share
\@namedef{val@tr05-p-boot}{0.033}% results/final/significance.csv
\@namedef{val@ld10-cascade}{0.8116}% results/final/metrics.csv (mean over 5 seeds)
\@namedef{val@ld10-svmchar}{0.8843}% results/final/metrics.csv (mean over 5 seeds)
\@namedef{val@ld01-cascade}{0.6957}% results/final/metrics.csv (mean over 5 seeds)
\@namedef{val@ld01-svmchar}{0.7073}% results/final/metrics.csv (mean over 5 seeds)
\@namedef{val@ld01-distilaug}{0.5418}% results/final/metrics.csv (mean over 5 seeds)
\@namedef{val@ld01-share}{17.9}% results/final/metrics.csv (mean over 5 seeds)
\@namedef{val@ld01-b}{62}% results/final/significance.csv
\@namedef{val@ld01-c}{18}% results/final/significance.csv
\@namedef{val@ld01-average}{0.7381}% results/final/metrics.csv (mean over 5 seeds)
\@namedef{val@ld05-average}{0.8476}% results/final/metrics.csv (mean over 5 seeds)
\@namedef{val@tr-zs-best}{0.6133}% results/final/metrics.csv (mean over 5 seeds)
\@namedef{val@tau-tr-zs}{0.9840}% results/final/tau.json
\@namedef{val@tau-tr-ft}{0.9985}% results/final/tau.json
\@namedef{val@tau-mh}{0.9990}% results/final/tau.json
\@namedef{val@tr-zs-share}{68.4}% results/final/metrics.csv (mean over 5 seeds)
\@namedef{val@tr-ft-share}{77.3}% results/final/metrics.csv (mean over 5 seeds)
\@namedef{val@tr-ft-cascade}{0.9705}% results/final/metrics.csv (mean over 5 seeds)
\@namedef{val@tr-ft-distilaug}{0.9645}% results/final/metrics.csv (mean over 5 seeds)
\@namedef{val@tr-ft-average}{0.9741}% results/final/metrics.csv (mean over 5 seeds)
\@namedef{val@tr-ft-p-boot}{0.034}% results/final/significance.csv
\@namedef{val@cost-clean-r1-cascade}{0.0064}% results/final/metrics.csv (mean over 5 seeds)
\@namedef{val@cost-clean-r3-cascade}{0.0378}% results/final/metrics.csv (mean over 5 seeds)
\@namedef{val@cost-clean-r3-distilaug}{0.0368}% results/final/metrics.csv (mean over 5 seeds)
\@namedef{val@cost-clean-r3-average}{0.0338}% results/final/metrics.csv (mean over 5 seeds)
\@namedef{val@cost-ho20-r2-cascade}{0.0260}% results/final/metrics.csv (mean over 5 seeds)
\@namedef{val@cost-ho20-r2-svmcharaug}{0.0303}% results/final/metrics.csv (mean over 5 seeds)
\@namedef{val@cost-ho20-r2-distilaug}{0.0364}% results/final/metrics.csv (mean over 5 seeds)
\@namedef{val@cost-mh-ho20-r2-cascade}{0.0301}% results/final/metrics.csv (mean over 5 seeds)
\@namedef{val@cost-mh-ho20-r2-svmcharaug}{0.0202}% results/final/metrics.csv (mean over 5 seeds)
\@namedef{val@lat-cascade-mean}{3.37}% results/final/cascade_end_to_end.json
\@namedef{val@lat-svmcharaug-thr}{669}% results/final/latency.csv
\@namedef{val@lat-bert-p50}{84.7}% results/final/latency.csv
\@namedef{val@lat-qwen-p50}{178}% results/final/latency.csv (GPU)
\@namedef{val@size-cascade}{273}% results/final/cascade_end_to_end.json
\@namedef{val@v2-slowdown}{48}% results/final/latency.csv / e2e p50
\@namedef{val@curve-clean-0}{0.9310}% results/final/cascade_curve.csv (mean over seeds)
\@namedef{val@curve-clean-peak}{0.9373}% results/final/cascade_curve.csv
\@namedef{val@curve-clean-peak-share}{1}% results/final/cascade_curve.csv
\@namedef{val@curve-clean-100}{0.9149}% results/final/cascade_curve.csv
\@namedef{val@noaug-wins}{6}% results/final/metrics.csv (mean over 5 seeds) cascade_noaug > cascade
\@namedef{val@n-rob-sets}{7}% configs/final.yaml mendeley_3class test_sets
\@namedef{val@noaug-beyond}{1}% results/final/metrics.csv (mean over 5 seeds) wins beyond the seed spread
\@namedef{val@men-smish-uci-spam}{242}% reports/stats/overlap.json crosstab (clean)
\@namedef{val@men-uci-smish}{495}% reports/stats/dedup_mendeley.json labels_after
\@namedef{val@infl-mh-acc}{0.24}% results/step5a/inflation (dirty - clean, pp)
\@namedef{val@infl-mh-f1}{0.96}% results/step5a/inflation (dirty - clean, pp)
\@namedef{val@infl-mh-acc-dirty}{0.9824}% results/step5a/inflation
\@namedef{val@infl-max-drop}{6.43}% results/step5a/inflation (largest macro F1 drop)
\@namedef{val@kz-orig}{42}% data/kz_cases/cases.csv
\@namedef{val@kz-all}{63}% data/kz_cases/cases.csv
\@namedef{val@kz-ru-missed-lite}{6}% results/step6/predictions.csv
\@namedef{val@kz-ru-missed-bert}{4}% results/step6/predictions.csv
\@namedef{val@kz-fraud-lang}{7}% results/step6/predictions.csv
\@namedef{val@kz-kk-missed-lite}{6}% results/step6/predictions.csv
\@namedef{val@kz-kk-missed-bert}{5}% results/step6/predictions.csv
\@namedef{val@kz-ru-tr-missed-min}{2}% results/step6/predictions.csv
\@namedef{val@kz-ru-tr-missed-max}{3}% results/step6/predictions.csv
\@namedef{val@kz-en-correct-lite}{7}% results/step6/predictions.csv
\@namedef{val@kz-en-n}{14}% results/step6/predictions.csv
\@namedef{val@kz-link-n}{25}% results/step6b/mt_sentence/translations.csv (RU and KK with a link)
\@namedef{val@kz-link-lost-word}{15}% results/step6/translations.csv
\@namedef{val@kz-ru-sent-before}{8}% results/step6/predictions.csv
\@namedef{val@kz-ru-sent-after}{11}% results/step6b/mt_sentence/predictions.csv
\@namedef{val@kz-ru-orig-n}{14}% results/step6b/mt_sentence/predictions.csv
\@namedef{val@kz-kk-sent-before}{10}% results/step6/predictions.csv
\@namedef{val@kz-kk-sent-after}{8}% results/step6b/mt_sentence/predictions.csv
\@namedef{val@kz-mt-p50}{4.6}% results/step6/latency.json
\@namedef{val@kz-mts-p50}{6.3}% results/step6b/mt_sentence/latency.json
\@namedef{val@kz-adapt-n}{324}% data/kz_adapt/adapt.csv
\@namedef{val@kz-adapt-ideas}{108}% data/kz_adapt/adapt.csv (ids without language)
\@namedef{val@kz-fraud-before}{7}% results/step6b/adapt/predictions.csv
\@namedef{val@kz-acc-before}{0.36}% results/step6b/adapt/predictions.csv
\@namedef{val@kz-fraud-after-a}{21}% results/step6b/adapt/predictions.csv
\@namedef{val@kz-acc-after-a}{0.74}% results/step6b/adapt/predictions.csv
\@namedef{val@kz-fraud-after-b}{21}% results/step6b/adapt/predictions.csv
\@namedef{val@kz-acc-after-b}{0.81}% results/step6b/adapt/predictions.csv
\@namedef{val@kz-fraud-n}{21}% results/step6b/adapt/predictions.csv
\@namedef{val@kz-val-min}{0.896}% results/step6b/adapt/runs.csv (mean per condition and size)
\@namedef{val@kz-val-max}{0.910}% results/step6b/adapt/runs.csv (mean per condition and size)
\@namedef{val@kz-val-before}{0.904}% results/step6b/adapt/runs.csv
\@namedef{val@kz-legit-link-min}{0.39}% results/step6b/adapt/predictions.csv (legitimate held-out messages with a link)
\@namedef{val@kz-legit-link-max}{0.62}% results/step6b/adapt/predictions.csv (legitimate held-out messages with a link)
\@namedef{val@kz-legit-link-before}{0.40}% results/step6b/adapt/predictions.csv
\@namedef{val@kz-legit-link-smish-before}{6}% results/step6b/adapt/predictions.csv
\@namedef{val@kz-legit-link-smish-after}{42}% results/step6b/adapt/predictions.csv (all adapted runs pooled)
\@namedef{val@kz-real-n}{4}% data/kz_real/public_smishing.csv
\@namedef{val@kz-real-flagged}{3}% results/step6c/predictions.csv (as written, adapted A and B)
\@namedef{val@kz-real-conf-min}{0.43}% results/step6c/predictions.csv
\@namedef{val@kz-real-conf-max}{0.66}% results/step6c/predictions.csv
\@namedef{val@kz-real-never}{1}% results/step6c/predictions.csv (no model, no input)
\@namedef{val@sh-clean-min}{0.8592}% results/final/metrics.csv (mean over 5 seeds)
\@namedef{val@sh-clean-max}{0.9027}% results/final/metrics.csv (mean over 5 seeds)
\@namedef{val@svmchar-shim-clean}{0.9276}% results/final/metrics.csv (mean over 5 seeds)
\@namedef{val@svmchar-mahmud-clean}{0.9864}% results/final/metrics.csv (mean over 5 seeds)
\@namedef{val@mh-clean}{0.9712}% results/final/metrics.csv (mean over 5 seeds)
\@namedef{val@lat-svmchar-p50}{1.47}% results/final/latency.csv
\@namedef{val@lat-svmchar-thr}{660}% results/final/latency.csv
\@namedef{val@lat-sh-p50-min}{43.2}% results/final/latency.csv (Shim et al., setting_shim)
\@namedef{val@lat-sh-p50-max}{89.5}% results/final/latency.csv (Shim et al., setting_shim)
\@namedef{val@speed-ratio-min}{29}% results/final/latency.csv p50 Shim / p50 svm_char
\@namedef{val@speed-ratio-max}{61}% results/final/latency.csv p50 Shim / p50 svm_char
\@namedef{val@size-svmchar}{5.4}% results/final/latency.csv size_mb
\@namedef{val@size-qwen-gb}{5.4}% results/final/latency.csv size_mb / 1000
\@namedef{val@explain-top}{15}% src/step8/explain.py TOP
\@namedef{val@ngrams-per-char}{14}% configs/step4c.yaml char_branch (n-grams covering one character)
\@namedef{val@ex-n-features}{136\,176}% results/step8/model.json n_features
\@namedef{val@ng-ham-1}{\texttt{..}}% results/step8/top_ngrams.csv
\@namedef{val@ng-ham-2}{\texttt{i\textvisiblespace{}}}% results/step8/top_ngrams.csv
\@namedef{val@ng-ham-3}{\texttt{\textvisiblespace{}i\textvisiblespace{}}}% results/step8/top_ngrams.csv
\@namedef{val@ng-ham-4}{\texttt{:)}}% results/step8/top_ngrams.csv
\@namedef{val@ng-ham-5}{\texttt{\textvisiblespace{}i}}% results/step8/top_ngrams.csv
\@namedef{val@ng-ham-6}{\texttt{my\textvisiblespace{}}}% results/step8/top_ngrams.csv
\@namedef{val@ng-ham-7}{\texttt{so}}% results/step8/top_ngrams.csv
\@namedef{val@ng-ham-8}{\texttt{\textvisiblespace{}m}}% results/step8/top_ngrams.csv
\@namedef{val@ng-ham-9}{\texttt{...}}% results/step8/top_ngrams.csv
\@namedef{val@ng-ham-10}{\texttt{me\textvisiblespace{}}}% results/step8/top_ngrams.csv
\@namedef{val@ng-ham-11}{\texttt{\textvisiblespace{}me}}% results/step8/top_ngrams.csv
\@namedef{val@ng-ham-12}{\texttt{\textvisiblespace{}my\textvisiblespace{}}}% results/step8/top_ngrams.csv
\@namedef{val@ng-ham-13}{\texttt{my}}% results/step8/top_ngrams.csv
\@namedef{val@ng-ham-14}{\texttt{\textvisiblespace{}me\textvisiblespace{}}}% results/step8/top_ngrams.csv
\@namedef{val@ng-ham-15}{\texttt{..\textvisiblespace{}}}% results/step8/top_ngrams.csv
\@namedef{val@ng-smishing-1}{\texttt{\textvisiblespace{}0}}% results/step8/top_ngrams.csv
\@namedef{val@ng-smishing-2}{\texttt{00}}% results/step8/top_ngrams.csv
\@namedef{val@ng-smishing-3}{\texttt{08}}% results/step8/top_ngrams.csv
\@namedef{val@ng-smishing-4}{\texttt{09}}% results/step8/top_ngrams.csv
\@namedef{val@ng-smishing-5}{\texttt{0\textvisiblespace{}}}% results/step8/top_ngrams.csv
\@namedef{val@ng-smishing-6}{\texttt{20}}% results/step8/top_ngrams.csv
\@namedef{val@ng-smishing-7}{\texttt{all\textvisiblespace{}0}}% results/step8/top_ngrams.csv
\@namedef{val@ng-smishing-8}{\texttt{ll\textvisiblespace{}0}}% results/step8/top_ngrams.csv
\@namedef{val@ng-smishing-9}{\texttt{00\textvisiblespace{}}}% results/step8/top_ngrams.csv
\@namedef{val@ng-smishing-10}{\texttt{l\textvisiblespace{}0}}% results/step8/top_ngrams.csv
\@namedef{val@ng-smishing-11}{\texttt{\textvisiblespace{}08}}% results/step8/top_ngrams.csv
\@namedef{val@ng-smishing-12}{\texttt{claim}}% results/step8/top_ngrams.csv
\@namedef{val@ng-smishing-13}{\texttt{laim}}% results/step8/top_ngrams.csv
\@namedef{val@ng-smishing-14}{\texttt{aim}}% results/step8/top_ngrams.csv
\@namedef{val@ng-smishing-15}{\texttt{clai}}% results/step8/top_ngrams.csv
\@namedef{val@ng-spam-1}{\texttt{xt}}% results/step8/top_ngrams.csv
\@namedef{val@ng-spam-2}{\texttt{\textvisiblespace{}8}}% results/step8/top_ngrams.csv
\@namedef{val@ng-spam-3}{\texttt{\textvisiblespace{}to\textvisiblespace{}8}}% results/step8/top_ngrams.csv
\@namedef{val@ng-spam-4}{\texttt{o\textvisiblespace{}8}}% results/step8/top_ngrams.csv
\@namedef{val@ng-spam-5}{\texttt{to\textvisiblespace{}8}}% results/step8/top_ngrams.csv
\@namedef{val@ng-spam-6}{\texttt{txt}}% results/step8/top_ngrams.csv
\@namedef{val@ng-spam-7}{\texttt{tone}}% results/step8/top_ngrams.csv
\@namedef{val@ng-spam-8}{\texttt{tx}}% results/step8/top_ngrams.csv
\@namedef{val@ng-spam-9}{\texttt{88}}% results/step8/top_ngrams.csv
\@namedef{val@ng-spam-10}{\texttt{50}}% results/step8/top_ngrams.csv
\@namedef{val@ng-spam-11}{\texttt{\textvisiblespace{}txt}}% results/step8/top_ngrams.csv
\@namedef{val@ng-spam-12}{\texttt{10}}% results/step8/top_ngrams.csv
\@namedef{val@ng-spam-13}{\texttt{ton}}% results/step8/top_ngrams.csv
\@namedef{val@ng-spam-14}{\texttt{\textvisiblespace{}tx}}% results/step8/top_ngrams.csv
\@namedef{val@ng-spam-15}{\texttt{ms}}% results/step8/top_ngrams.csv
\@namedef{val@ngm-ham-1}{\texttt{..}}% results/step8/top_ngrams_long.csv (merged, make_tables.merged_ngrams)
\@namedef{val@ngm-ham-2}{\texttt{i\textvisiblespace{}}}% results/step8/top_ngrams_long.csv (merged, make_tables.merged_ngrams)
\@namedef{val@ngm-ham-3}{\texttt{:)}}% results/step8/top_ngrams_long.csv (merged, make_tables.merged_ngrams)
\@namedef{val@ngm-ham-4}{\texttt{my\textvisiblespace{}}}% results/step8/top_ngrams_long.csv (merged, make_tables.merged_ngrams)
\@namedef{val@ngm-ham-5}{\texttt{so}}% results/step8/top_ngrams_long.csv (merged, make_tables.merged_ngrams)
\@namedef{val@ngm-ham-6}{\texttt{\textvisiblespace{}m}}% results/step8/top_ngrams_long.csv (merged, make_tables.merged_ngrams)
\@namedef{val@ngm-ham-7}{\texttt{me\textvisiblespace{}}}% results/step8/top_ngrams_long.csv (merged, make_tables.merged_ngrams)
\@namedef{val@ngm-ham-8}{\texttt{\textvisiblespace{}it}}% results/step8/top_ngrams_long.csv (merged, make_tables.merged_ngrams)
\@namedef{val@ngm-smishing-1}{\texttt{\textvisiblespace{}0}}% results/step8/top_ngrams_long.csv (merged, make_tables.merged_ngrams)
\@namedef{val@ngm-smishing-2}{\texttt{00}}% results/step8/top_ngrams_long.csv (merged, make_tables.merged_ngrams)
\@namedef{val@ngm-smishing-3}{\texttt{08}}% results/step8/top_ngrams_long.csv (merged, make_tables.merged_ngrams)
\@namedef{val@ngm-smishing-4}{\texttt{09}}% results/step8/top_ngrams_long.csv (merged, make_tables.merged_ngrams)
\@namedef{val@ngm-smishing-5}{\texttt{20}}% results/step8/top_ngrams_long.csv (merged, make_tables.merged_ngrams)
\@namedef{val@ngm-smishing-6}{\texttt{all\textvisiblespace{}0}}% results/step8/top_ngrams_long.csv (merged, make_tables.merged_ngrams)
\@namedef{val@ngm-smishing-7}{\texttt{claim}}% results/step8/top_ngrams_long.csv (merged, make_tables.merged_ngrams)
\@namedef{val@ngm-smishing-8}{\texttt{50}}% results/step8/top_ngrams_long.csv (merged, make_tables.merged_ngrams)
\@namedef{val@ngm-spam-1}{\texttt{txt}}% results/step8/top_ngrams_long.csv (merged, make_tables.merged_ngrams)
\@namedef{val@ngm-spam-2}{\texttt{\textvisiblespace{}8}}% results/step8/top_ngrams_long.csv (merged, make_tables.merged_ngrams)
\@namedef{val@ngm-spam-3}{\texttt{\textvisiblespace{}to\textvisiblespace{}8}}% results/step8/top_ngrams_long.csv (merged, make_tables.merged_ngrams)
\@namedef{val@ngm-spam-4}{\texttt{tone}}% results/step8/top_ngrams_long.csv (merged, make_tables.merged_ngrams)
\@namedef{val@ngm-spam-5}{\texttt{88}}% results/step8/top_ngrams_long.csv (merged, make_tables.merged_ngrams)
\@namedef{val@ngm-spam-6}{\texttt{50}}% results/step8/top_ngrams_long.csv (merged, make_tables.merged_ngrams)
\@namedef{val@ngm-spam-7}{\texttt{10}}% results/step8/top_ngrams_long.csv (merged, make_tables.merged_ngrams)
\@namedef{val@ngm-spam-8}{\texttt{ms}}% results/step8/top_ngrams_long.csv (merged, make_tables.merged_ngrams)
\@namedef{val@ngm-smishing-digits}{6}% results/step8/top_ngrams_long.csv (merged smishing features of digits only)
\@namedef{val@ngm-top}{8}% tools/make_tables.py merged_ngrams top
\@namedef{val@ngm-pool}{60}% src/step8/explain.py TOP_LONG
\@namedef{val@ex-icpt}{$-$0.09}% results/step8/model.json intercepts (smishing - spam)
\@namedef{val@ex-top-words}{8}% tools/make_figures.py explanation_figure top
\@namedef{val@ex-sm-n}{91}% results/step8/examples.json n_in_group
\@namedef{val@ex-sm-conf}{0.97}% results/step8/examples.json
\@namedef{val@ex-sm-margin}{$+$1.53}% results/step8/examples.json scores (s_smishing - s_spam)
\@namedef{val@ex-sm-rest}{$+$0.29}% results/step8/example_words.csv (words outside the top 8)
\@namedef{val@ex-sm-nwords}{18}% results/step8/example_words.csv
\@namedef{val@ex-sm-url}{0.68}% results/step8/example_words.csv
\@namedef{val@ex-sm-second}{0.13}% results/step8/example_words.csv (second largest word)
\@namedef{val@ex-sp-n}{75}% results/step8/examples.json n_in_group
\@namedef{val@ex-sp-conf}{0.96}% results/step8/examples.json
\@namedef{val@ex-sp-margin}{$-$0.96}% results/step8/examples.json scores (s_smishing - s_spam)
\@namedef{val@ex-sp-rest}{$-$0.06}% results/step8/example_words.csv (words outside the top 8)
\@namedef{val@ex-sp-nwords}{24}% results/step8/example_words.csv
\@namedef{val@ex-sp-max}{0.15}% results/step8/example_words.csv (largest word)
\@namedef{val@ex-err-n}{7}% results/step8/examples.json n_in_group
\@namedef{val@ex-err-conf}{0.85}% results/step8/examples.json
\@namedef{val@ex-err-margin}{$-$0.93}% results/step8/examples.json scores (s_smishing - s_spam)
\@namedef{val@ex-err-rest}{$+$0.09}% results/step8/example_words.csv (words outside the top 8)
\@namedef{val@ex-err-nwords}{15}% results/step8/example_words.csv
\@namedef{val@ex-err-stop}{0.53}% results/step8/example_words.csv
\@namedef{val@ex-err-phone}{0.31}% results/step8/example_words.csv
\makeatother
\newcommand{\val}[1]{\ifcsname val@#1\endcsname\csname val@#1\endcsname\else\textbf{??#1??}\fi}


%=================================================================
\Title{Are Transformers Necessary for Smishing Detection? A Leakage-Free, Obfuscation-Aware Evaluation}

% Author Orchid ID: enter ID or remove command
\newcommand{\orcidauthorA}{0000-0000-0000-000X}

\Author{Firstname Lastname $^{1}$\orcidA{} and Firstname Lastname $^{2,}$*}

\AuthorNames{Firstname Lastname and Firstname Lastname}

\address{%
$^{1}$ \quad Affiliation 1; e-mail@e-mail.com\\
$^{2}$ \quad Affiliation 2; e-mail@e-mail.com}

\corres{Correspondence: e-mail@e-mail.com}

% Abstract (one paragraph: background, methods, results, conclusions; at most 200 words, checked by check_numbers.py)
\abstract{Fraudulent text messages (smishing) are often disguised with character substitutions, homoglyphs and inserted spaces, yet detectors are usually evaluated on clean corpora with duplicates. We present a leakage-free, obfuscation-aware evaluation of real-time smishing detectors on three deduplicated public corpora. Obfuscation is modelled as a random string transformation, with one family of operations for training and another held out for testing. We compare two operating points, Lite, a calibrated linear model over character n-grams for devices, and Full, a cascade that passes low-confidence messages to DistilBERT for operator gateways, with classical, neural and re-implemented published detectors (five seeds; paired bootstrap and McNemar tests with Holm correction). Under the strongest held-out obfuscation, the published transformer detectors fell from \val{sh-clean-min}--\val{sh-clean-max} to \val{sh-ho20-min}--\val{sh-ho20-max} macro F1, while Lite went from \val{svmchar-shim-clean} to \val{svmchar-shim-ho20} (differences significant by both tests). Lite took a median of \val{lat-svmchar-p50}~ms per message on one CPU thread, \val{speed-ratio-min} to \val{speed-ratio-max} times less than these transformers. Full sent \val{lat-cascade-routed-pct}\% of messages to DistilBERT and had the lowest error cost under strong obfuscation, but on clean data did not differ significantly from Lite (\val{cascade-clean-f1} versus \val{svmchar-clean-f1}). Calibrated character-level models are thus a strong default for real-time smishing detection.}

% Keywords
\keyword{smishing detection; fraudulent text messages; character n-grams; obfuscation robustness; probability calibration; cascade classification; data leakage; real-time inference}

%=================================================================
\begin{document}

\section{Introduction}\label{sec:intro}

Smishing, or phishing by short text message, imitates a bank, a delivery service or a government agency and leads the recipient to a link, a phone number or an application in order to obtain credentials or a payment. Unlike advertising spam, a missed smishing message can cause direct financial loss, so the two kinds of error have unequal costs. A detector for this task works under constraints that differ from those of document classification. The messages are short, informal and often written in several languages. Senders alter spelling, insert spaces or replace letters with visually similar symbols in order to evade filters \cite{salman2024,boucher2022}. A filter that runs on a handset or at an operator gateway must classify every message within milliseconds, and the content of private messages should not leave the device or the operator network.

Recent detectors rely more and more on fine-tuned transformer encoders \cite{ghourabi2023,shim2024} and on large language models, and many studies report high accuracy on public corpora of short messages \cite{mishra2023dsmishsms,mishra2022detector,sonowal2020,roy2020,mahmud2024}. Three aspects of the usual evaluation make it hard to tell whether this added capacity is needed in practice.

\emph{Duplicates and leakage.} The Mendeley SMS Phishing Dataset \cite{mendeley2022,mishra2023dataset} is one of the few public corpora that separates smishing from spam. Of its \val{men-raw} messages, our deduplication removed \val{men-removed}: exact or near duplicates of a kept message and members of \val{men-conflict-groups} groups of identical or nearly identical texts that carry different labels. Of the \val{men-clean} remaining messages, \val{men-in-uci} (\val{men-in-uci-pct}\%) have an exact or near duplicate in the older UCI SMS Spam Collection \cite{almeida2011}. A random split of such data can place copies of one message in both the training and the test part, a known source of leakage and of optimistic estimates \cite{kapoor2023,lee2022}. The published methods that we re-implement do not state whether duplicates were removed \cite{mahmud2024,shim2024}.

\emph{Obfuscation that the model has not seen.} Character-level perturbations degrade word-level and subword neural models \cite{gao2018,li2019,pruthi2019}, and Unicode look-alike characters can change the input of a model without changing what a reader sees \cite{boucher2022}. When robustness is measured at all, the test perturbations are usually of the same kind as those used in training, which says little about an attacker who uses a different trick.

\emph{Inference cost.} Accuracy is rarely reported together with latency, and almost never end to end for a complete detection pipeline. In our measurements on one CPU thread, a fine-tuned DistilBERT classifier needs a median of \val{lat-distilaug-p50}~ms per message, while a linear model over character n-grams needs \val{lat-svmchar-p50}~ms.

This paper asks whether transformers are necessary for real-time smishing detection once these three aspects are controlled. We evaluate detectors on three public corpora in English and Portuguese after removing exact and near duplicates, model obfuscation as a random string transformation with one family of operations for training and another family held out for testing, and measure latency on one CPU thread. The detector we study has two operating points. \emph{Lite} is a linear support vector machine over character n-grams with calibrated class probabilities, small and fast enough to run on a handset. \emph{Full} is a cascade for an operator gateway: it answers with the character model when its calibrated confidence reaches a threshold and passes the remaining messages to a DistilBERT classifier. Both are compared with classical, neural and transformer baselines, with a few-shot large language model \cite{qwen2024}, and with re-implementations of two recent published methods \cite{mahmud2024,shim2024} in their own task settings, under a protocol that was frozen before the test data were read.

The contributions are as follows.
\begin{enumerate}
\item \emph{A leakage-free, obfuscation-aware evaluation protocol.} It uses deduplicated splits of three corpora with conflicting labels resolved and no exact or near duplicates across the splits; one obfuscation family for training and a different one held out for testing; a configuration frozen before testing; \val{n-seeds} training seeds; paired bootstrap and McNemar tests with Holm correction; and re-implementation of published methods in the task settings of their papers. We also show that the treatment of duplicates alone changes cross-validated macro F1 by up to \val{infl-max-drop} points, more than most differences between models.
\item \emph{Evidence that published transformer detectors are fragile under unseen obfuscation.} In the smishing-versus-spam setting of Shim et al. \cite{shim2024}, the re-implemented BERT, DistilBERT and RoBERTa detectors reach between \val{sh-clean-min} and \val{sh-clean-max} macro F1 on clean test data and fall to between \val{sh-ho20-min} and \val{sh-ho20-max} when \val{lvl-c}\% of the characters are edited with the held-out operations. Lite goes from \val{svmchar-shim-clean} to \val{svmchar-shim-ho20} on the same data, and its advantage at the \val{lvl-b}\% and \val{lvl-c}\% levels is significant by both tests. Lite needs \val{lat-svmchar-p50}~ms per message, \val{speed-ratio-min} to \val{speed-ratio-max} times less than these transformers.
\item \emph{Two operating points with an explicit choice of the threshold.} The threshold of Full is the solution of a stated selection problem on validation data. In the three-class task, Full sends \val{lat-cascade-routed-pct}\% of the test messages to DistilBERT, has a median end-to-end latency of \val{lat-cascade-p50}~ms and the lowest error cost of all models under the strongest held-out obfuscation. On clean data, Full does not differ significantly from Lite.
\item \emph{Interpretation and deployment.} We read the n-gram weights of the trained Lite model to show which character patterns drive its decisions, present a case study for Kazakhstan with messages in Russian and Kazakh, and derive practical recommendations for banks and mobile operators.
\end{enumerate}

Section~\ref{sec:related} reviews related work. Section~\ref{sec:threat} states the threat model. Section~\ref{sec:methods} describes the data, the obfuscation model, Lite and Full, the cost function, the statistical tests and the protocol. Section~\ref{sec:results} presents the results, Section~\ref{sec:discussion} discusses them, and Section~\ref{sec:conclusions} concludes.

\section{Related Work}\label{sec:related}

[Section in preparation.]

\section{Threat Model}\label{sec:threat}

This section states what the attacker can do, what constrains the defender, and which part of this threat our evaluation covers. It describes the setting of the study and involves no experiment.

\subsection{Attacker}\label{sec:threat-attacker}

The attacker sends text messages that imitate a trusted sender and wants them to reach the recipient unflagged. A message has to stay readable and persuasive for a person, so every evasion is limited by what a reader still understands. Within this limit, the attacker can:
\begin{itemize}
\item replace characters by look-alike ASCII symbols (0 for o, 1 for l, @ for a, \$ for s);
\item replace Latin letters by visually identical letters of another script, such as Cyrillic, or insert invisible Unicode characters \cite{boucher2022};
\item insert spaces inside words, double letters or drop vowels, which keeps a word recognisable while changing its spelling;
\item paraphrase the message, by hand or with a language model, so that the wording differs from known campaigns while the request stays the same;
\item write in another language or script, for example in Russian or Kazakh, or transliterate Cyrillic text into Latin letters;
\item adapt to the detector: send variants, observe which ones are delivered and keep the edits that pass (black-box access), or, in the worst case, compute edits from the model itself (white-box access).
\end{itemize}
The attacker's goal in this study is evasion at test time. Poisoning of the training data \cite{tang2025} and attacks on the link target, such as URL shorteners or compromised domains, are outside the scope; our detectors see only the text of a message.

\subsection{Defender}\label{sec:threat-defender}

The defender is either a handset application or a filter at the gateway of a mobile operator, and both work under constraints that rule out some detectors.
\begin{itemize}
\item \emph{Latency and throughput.} A gateway must classify every message as it passes, and a handset must not delay the notification noticeably. The cost per message, not only the accuracy, decides whether a detector can be deployed.
\item \emph{Memory and energy.} On a handset the model has to be small and cheap to evaluate on a CPU, without a GPU.
\item \emph{Privacy.} Private messages should not be sent to an external service for classification. A detector that runs on the device or inside the operator network keeps the content local, while a hosted large language model does not.
\item \emph{Unequal error costs.} A missed smishing message can cause direct financial loss, while a false alarm costs the attention of the user or a delayed legitimate notification. We express this by the cost ratio $r$ in Section~\ref{sec:cost}.
\end{itemize}
These constraints lead to the two operating points of this study: Lite for a handset, where only a small model fits, and Full for a gateway, where a transformer can be afforded for a small share of the messages.

\subsection{What the Evaluation Covers}\label{sec:threat-scope}

Table~\ref{tab:threat} maps the attacker capabilities to the evaluation. Obfuscation is modelled as a random string transformation (Section~\ref{sec:obfuscation}) with two families of operations. The family TR (look-alike substitution and inserted spaces) is used for training augmentation and for testing. The family HO (Latin-to-Cyrillic homoglyphs, letter doubling and vowel deletion) is never used in training and models an attacker whose trick the defender has not anticipated. Each family is tested at \val{lvl-a}\%, \val{lvl-b}\% and \val{lvl-c}\% of the characters of a message.

\begin{table}[H]
\caption{Attacker capabilities and how they are covered by the evaluation. TR: obfuscation family used in training; HO: held-out family, used only at test time.\label{tab:threat}}
\begin{tabularx}{\textwidth}{L L L}
\toprule
\textbf{Capability} & \textbf{Modelled} & \textbf{Where evaluated}\\
\midrule
Look-alike ASCII substitution & yes, family TR & augmentation; TR test sets\\
Inserted spaces & yes, family TR & augmentation; TR test sets\\
Latin-to-Cyrillic homoglyphs & yes, family HO & HO test sets only\\
Letter doubling, vowel deletion & yes, family HO & HO test sets only\\
Cyrillic-to-Latin look-alikes & partly & Kazakhstan case study (qualitative)\\
Another language (Russian, Kazakh) & partly & Kazakhstan case study (qualitative); transfer to Portuguese\\
Invisible characters, reordering & no & \\
Paraphrasing, generated rewording & no & \\
Adaptive attacks (black-box or white-box) & no & \\
Data poisoning, attacks on the link target & no (out of scope) & \\
\bottomrule
\end{tabularx}
\end{table}

The obfuscation model is random and non-adaptive: positions and operations are drawn independently of the detector, the edits do not interact, and the attacker has no access to the model. An attacker who chooses edits to fool a particular detector, or who paraphrases the message, can do at least as well. Our results under obfuscation are therefore an optimistic estimate of the robustness of every detector, including ours, and their main value is the comparison between detectors under the same perturbations. Since the homoglyph and vowel operations act on Latin letters only, on Cyrillic text the held-out family reduces mostly to letter doubling.

\section{Materials and Methods}\label{sec:methods}

[Section in preparation.]

\section{Results}\label{sec:results}

[Section in preparation.]

\section{Discussion}\label{sec:discussion}

[Section in preparation.]

\section{Conclusions}\label{sec:conclusions}

[Section in preparation.]


\vspace{6pt}

%% Back matter in the order of the MDPI template. Text in square brackets is a placeholder for the authors.

\authorcontributions{[To be completed by the authors, following the CRediT taxonomy: conceptualization, methodology, software, validation, formal analysis, investigation, data curation, writing (original draft preparation), writing (review and editing), visualization, supervision.] All authors have read and agreed to the published version of the manuscript.}

\funding{[To be completed by the authors: ``This research received no external funding'' or the funder and grant number.]}

\institutionalreview{Not applicable.}

\informedconsent{Not applicable.}

% left-aligned: the long URLs would leave wide gaps in a justified paragraph (\par of the class runs inside this group)
{\raggedright
\dataavailability{The corpora analysed in this study are public and are not redistributed with this article: the Mendeley SMS Phishing Dataset (\url{https://doi.org/10.17632/f45bkkt8pr.1}, licence CC BY 4.0) \cite{mendeley2022}, MOZ-Smishing (\url{https://huggingface.co/datasets/MOZNLP/MOZ-Smishing}) \cite{moz2025}, the Super SMS Dataset (\url{https://github.com/smspamresearch/spstudy}; no licence is stated for the data) \cite{salman2024} and the UCI SMS Spam Collection (\url{https://archive.ics.uci.edu/dataset/228/sms+spam+collection}) \cite{almeida2011}. The repository of this study contains the scripts that download these files, check their checksums and rebuild the deduplicated splits, the identifiers of the messages in every split, the lists of removed duplicates and conflicting groups, all predictions and result files, the author-written case-study messages and the synthetic adaptation set. The four real messages of the case study are quoted from public news articles listed in the repository. All code for data preparation, training, evaluation, statistical tests, tables and figures, and the frozen configuration of the final evaluation, is in the same repository; every number in this article is generated from the result files by a script, and a second script checks that the text and tables match these files. Repository: [URL to be added by the authors].}
}

\acknowledgments{[Acknowledgments of support, if any, to be completed by the authors.] During the preparation of this work, the authors used Claude (Anthropic; model Claude Opus 5.5) to assist with code development, data processing scripts, generation of the synthetic adaptation set, and language editing of the manuscript. The authors reviewed and verified all generated content and take full responsibility for the content of this publication.}

\conflictsofinterest{[To be completed by the authors: ``The authors declare no conflicts of interest'' or the relevant declaration.]}

\abbreviations{Abbreviations}{%
The following abbreviations are used in this manuscript:\\

\noindent
\begin{tabular}{@{}ll}
CPU & Central processing unit\\
ECE & Expected calibration error\\
GPU & Graphics processing unit\\
HO & Held-out obfuscation family\\
LLM & Large language model\\
SMS & Short message service\\
SVM & Support vector machine\\
TF-IDF & Term frequency, inverse document frequency\\
TR & Obfuscation family used in training
\end{tabular}
}


\appendixstart
\appendix
\section[\appendixname~\thesection]{Additional Tables and Figures}\label{app:full}

[Appendix in preparation.]


%=================================================================
\begin{adjustwidth}{-\extralength}{0cm}

\reftitle{References}
% Generated by paper_mdpi/tools/make_bibliography.py from references.bib (citation order). Do not edit.
\begin{thebibliography}{999}

\bibitem{salman2024}
Salman, M.; Ikram, M.; Kaafar, M.A. Investigating Evasive Techniques in SMS Spam Filtering: A Comparative Analysis of Machine Learning Models. \emph{IEEE Access} \textbf{2024}, \emph{12}, 24306--24324. \url{https://doi.org/10.1109/access.2024.3364671}.

\bibitem{boucher2022}
Boucher, N.; Shumailov, I.; Anderson, R.; Papernot, N. Bad Characters: Imperceptible NLP Attacks. In \emph{2022 IEEE Symposium on Security and Privacy (SP)}; IEEE, 2022; pp. 1987--2004. \url{https://doi.org/10.1109/sp46214.2022.9833641}.

\bibitem{ghourabi2023}
Ghourabi, A.; Alohaly, M. Enhancing Spam Message Classification and Detection Using Transformer-Based Embedding and Ensemble Learning. \emph{Sensors} \textbf{2023}, \emph{23}, 3861. \url{https://doi.org/10.3390/s23083861}.

\bibitem{shim2024}
Shim, H.S.; Park, H.; Lee, K.; Park, J.-S.; Kang, S. A Persuasion-Based Prompt Learning Approach to Improve Smishing Detection through Data Augmentation. \emph{arXiv} \textbf{2024}, arXiv:2411.02403.

\bibitem{mishra2023dsmishsms}
Mishra, S.; Soni, D. DSmishSMS-A System to Detect Smishing SMS. \emph{Neural Computing and Applications} \textbf{2023}, \emph{35}, 4975--4992. \url{https://doi.org/10.1007/s00521-021-06305-y}.

\bibitem{mishra2022detector}
Mishra, S.; Soni, D. Implementation of `Smishing Detector': An Efficient Model for Smishing Detection Using Neural Network. \emph{SN Computer Science} \textbf{2022}, \emph{3}, 189. \url{https://doi.org/10.1007/s42979-022-01078-0}.

\bibitem{sonowal2020}
Sonowal, G. Detecting Phishing SMS Based on Multiple Correlation Algorithms. \emph{SN Computer Science} \textbf{2020}, \emph{1}, 361. \url{https://doi.org/10.1007/s42979-020-00377-8}.

\bibitem{roy2020}
Roy, P.K.; Singh, J.P.; Banerjee, S. Deep learning to filter SMS Spam. \emph{Future Generation Computer Systems} \textbf{2020}, \emph{102}, 524--533. \url{https://doi.org/10.1016/j.future.2019.09.001}.

\bibitem{mahmud2024}
Mahmud, T.; Prince, M.A.H.; Ali, M.H.; Hossain, M.S.; Andersson, K. Enhancing Cybersecurity: Hybrid Deep Learning Approaches to Smishing Attack Detection. \emph{Systems} \textbf{2024}, \emph{12}, 490. \url{https://doi.org/10.3390/systems12110490}.

\bibitem{mendeley2022}
Mishra, S.; Soni, D. SMS Phishing Dataset for Machine Learning and Pattern Recognition. Mendeley Data, version 1, 2022. \url{https://doi.org/10.17632/f45bkkt8pr.1}.

\bibitem{mishra2023dataset}
Mishra, S.; Soni, D. SMS Phishing Dataset for Machine Learning and Pattern Recognition. In \emph{Proceedings of the 14th International Conference on Soft Computing and Pattern Recognition (SoCPaR 2022)}; Lecture Notes in Networks and Systems; Springer, Cham, 2023; pp. 597--604. \url{https://doi.org/10.1007/978-3-031-27524-1_57}.

\bibitem{almeida2011}
Almeida, T.A.; Hidalgo, J.M.G.; Yamakami, A. Contributions to the study of SMS spam filtering. In \emph{Proceedings of the 11th ACM symposium on Document engineering}; ACM, 2011; pp. 259--262. \url{https://doi.org/10.1145/2034691.2034742}.

\bibitem{kapoor2023}
Kapoor, S.; Narayanan, A. Leakage and the reproducibility crisis in machine-learning-based science. \emph{Patterns} \textbf{2023}, \emph{4}, 100804. \url{https://doi.org/10.1016/j.patter.2023.100804}.

\bibitem{lee2022}
Lee, K.; Ippolito, D.; Nystrom, A.; Zhang, C.; Eck, D.; Callison-Burch, C.; Carlini, N. Deduplicating Training Data Makes Language Models Better. In \emph{Proceedings of the 60th Annual Meeting of the Association for Computational Linguistics (Volume 1: Long Papers)}; Association for Computational Linguistics, 2022; pp. 8424--8445. \url{https://doi.org/10.18653/v1/2022.acl-long.577}.

\bibitem{gao2018}
Gao, J.; Lanchantin, J.; Soffa, M.L.; Qi, Y. Black-Box Generation of Adversarial Text Sequences to Evade Deep Learning Classifiers. In \emph{2018 IEEE Security and Privacy Workshops (SPW)}; IEEE, 2018; pp. 50--56. \url{https://doi.org/10.1109/spw.2018.00016}.

\bibitem{li2019}
Li, J.; Ji, S.; Du, T.; Li, B.; Wang, T. TextBugger: Generating Adversarial Text Against Real-world Applications. In \emph{Proceedings 2019 Network and Distributed System Security Symposium}; Internet Society, 2019. \url{https://doi.org/10.14722/ndss.2019.23138}.

\bibitem{pruthi2019}
Pruthi, D.; Dhingra, B.; Lipton, Z.C. Combating Adversarial Misspellings with Robust Word Recognition. In \emph{Proceedings of the 57th Annual Meeting of the Association for Computational Linguistics}; Association for Computational Linguistics, 2019; pp. 5582--5591. \url{https://doi.org/10.18653/v1/p19-1561}.

\bibitem{qwen2024}
Qwen Team Qwen2.5 Technical Report. \emph{arXiv} \textbf{2024}, arXiv:2412.15115.

\bibitem{moz2025}
Ali, F.D.M.A.; Saide, S.M.; Sousa-Silva, R.; Lopes Cardoso, H. MOZ-Smishing: A Benchmark Dataset for Detecting Mobile Money Frauds. In \emph{Proceedings of the Sixth Workshop on African Natural Language Processing (AfricaNLP 2025)}; Association for Computational Linguistics, 2025; pp. 158--166. \url{https://doi.org/10.18653/v1/2025.africanlp-1.23}.

\end{thebibliography}


\PublishersNote{}
\end{adjustwidth}
\end{document}
